% !TEX root = ../main.tex
\subsection{Compositional structure of the reduced generator}
\label{subsec:compositional_structure}

The Kolmogorov--Arnold theorem asks when a multivariate map can be written as a
finite composition of univariate functions and addition; in the notation used
by Kolmogorov--Arnold networks one seeks representations of the form
\(F(\mathbf x)=\sum_q \Phi_q(\sum_p \phi_{q,p}(x_p))\), with the univariate
edge functions then learned inside a chosen basis
class~\cite{liuKANKolmogorovArnoldNetworks2024}. The relevant point for the
present paper is that the exact Gaussian HMF already comes with its own
compositional architecture. The KAN comparison therefore does not need to be
metaphorical.

\begin{proposition}
For a Gaussian bath linearly coupled through $f$, the exact Hamiltonian of mean
force is generated by three explicit compositional layers determined entirely
by $(H_Q,f,K,\beta)$. First, the univariate operator orbit is the adjoint chain
\(f_n=\mathrm{ad}_{H_Q}^n(f)\). Second, the influence exponent is the bilinear
aggregate
\begin{equation}
    \Delta(\beta)=\sum_{n,m\ge0} C^>_{nm}(\beta)\,f_n f_m .
\end{equation}
Third, the reduced generator is obtained by nonlinear BCH recombination of that
aggregate with the bare Gibbs factor,
\begin{equation}
    H_{\mathrm{MF}}(\beta)
    = H_Q-\frac{1}{\beta}\Phi\!\big(-\beta\,\mathrm{ad}_{H_Q}\big)\Delta(\beta)
    + O(\Delta^2),
\end{equation}
with the higher-order terms supplied by the full nested BCH series. Conditions
(C1)--(C3) of the closure theorem are exactly the layerwise conditions that
keep the adjoint orbit, the bilinear aggregate, and the nonlinear BCH tower
inside a chosen operator ansatz $\mathcal A$.
\end{proposition}

Figure~\ref{fig:kan_hmf_closure} makes this three-layer correspondence explicit
and shows the spin-boson qubit as the worked case in which each open-ended
layer terminates into a closed finite structure.

\begin{figure*}[t]
    \centering
    \resizebox{0.96\textwidth}{!}{\begin{tikzpicture}[
    x=1cm,
    y=1cm,
    >=Latex,
    font=\small,
    line cap=round,
    line join=round,
    header/.style={font=\bfseries\small, align=center},
    layerlab/.style={font=\scriptsize, text=black!68, anchor=east},
    phimark/.style={font=\normalsize\bfseries, text=accent},
    plainarrow/.style={-{Latex[length=2mm]}, thin},
    corresp/.style={densely dashed, -{Latex[length=2mm]}, thin, draw=black!58},
    orbitnode/.style={
        draw,
        rounded corners=1pt,
        inner sep=2pt,
        minimum width=1.05cm,
        minimum height=0.55cm,
        fill=white
    },
    box/.style={
        draw,
        rounded corners=2pt,
        thin,
        fill=white,
        inner sep=4pt,
        align=center
    },
    smalllab/.style={font=\scriptsize, inner sep=1pt}
]
\colorlet{accent}{blue!70!black}
\colorlet{muted}{black!68}

\def\xL{2.55}
\def\xC{8.55}
\def\xR{14.65}
\def\yIn{0.70}
\def\yAtoms{2.15}
\def\yAgg{5.05}
\def\yOuter{7.78}
\def\yOut{10.05}
\def\yHead{11.25}

% Headers
\node[header] at (\xL,\yHead) {Kolmogorov--Arnold\\representation};
\node[header] at (\xC,\yHead) {Mean-force\\representation};
\node[header] at (\xR,\yHead) {Spin-boson qubit\\(exact closure)};

% Stratum labels
\node[layerlab] at (0.15,\yAtoms) {Atoms};
\node[layerlab] at (0.15,\yAgg) {Aggregation};
\node[layerlab] at (0.15,\yOuter) {Outer composition};

% Horizontal correspondence arrows
\draw[corresp] (4.85,\yAtoms+0.15) -- (6.45,\yAtoms+0.15);
\draw[corresp] (4.85,\yAgg) -- (6.45,\yAgg);
\draw[corresp] (4.85,\yOuter) -- (6.45,\yOuter);
\draw[corresp] (11.20,\yAtoms+0.15) -- (13.00,\yAtoms+0.15);
\draw[corresp] (11.20,\yAgg) -- (13.00,\yAgg);
\draw[corresp] (11.20,\yOuter) -- (13.00,\yOuter);
\node[font=\scriptsize\bfseries, text=muted] at (12.15,\yAtoms+0.47) {$C_1$};
\node[font=\scriptsize\bfseries, text=muted] at (12.15,\yAgg+0.33) {$C_2$};
\node[font=\scriptsize\bfseries, text=muted] at (12.15,\yOuter+0.33) {$C_3$};

% Left column: KAN inputs
\node[box, text width=2.20cm] (outL) at (\xL,\yOut) {$f(\mathbf{x})$};
\node[smalllab, text=muted] at (\xL+1.88,8.90) {$\vdots$};
\node[smalllab, text=muted] at (\xL+1.88,6.70) {$\vdots$};

\coordinate (LPhi) at (\xL,\yOuter);
\draw[accent, line width=1.05pt]
    ($(LPhi)+(-0.80,-0.12)$)
    .. controls +(0.35,0.55) and +(-0.40,-0.25) ..
    ($(LPhi)+(-0.14,0.42)$)
    .. controls +(0.38,0.28) and +(-0.24,-0.38) ..
    ($(LPhi)+(0.30,0.03)$)
    .. controls +(0.20,0.32) and +(-0.24,-0.10) ..
    ($(LPhi)+(0.80,0.30)$);
\draw[accent, line width=0.8pt]
    ($(LPhi)+(-0.74,-0.34)$)
    .. controls +(0.28,0.20) and +(-0.22,-0.24) ..
    ($(LPhi)+(-0.08,-0.08)$)
    .. controls +(0.28,0.18) and +(-0.22,-0.16) ..
    ($(LPhi)+(0.60,-0.24)$);
\node[phimark] at ($(LPhi)+(0,-0.67)$) {$\Phi_q$};
\node[smalllab, text=muted] at ($(LPhi)+(0,0.73)$) {learned};
\draw[plainarrow] ($(LPhi)+(0,0.62)$) -- (outL.south);
\node[smalllab, text=muted] at (\xL+0.92,8.55) {$\sum_q$};

\node[box, text width=3.10cm] (LAgg) at (\xL,\yAgg)
    {$\displaystyle \sum_p \phi_{q,p}(x_p)$};
\draw[plainarrow] (LAgg.north) -- ($(LPhi)+(0,-0.22)$);

\coordinate (x1) at (0.95,\yIn);
\coordinate (x2) at (1.85,\yIn);
\coordinate (x3) at (2.75,\yIn);
\coordinate (xn) at (4.15,\yIn);
\coordinate (p1) at (0.95,\yAtoms);
\coordinate (p2) at (1.85,\yAtoms);
\coordinate (p3) at (2.75,\yAtoms);
\coordinate (pn) at (4.15,\yAtoms);

\node[smalllab] at (x1) {$x_1$};
\node[smalllab] at (x2) {$x_2$};
\node[smalllab] at (x3) {$x_3$};
\node[smalllab] at (3.38,\yIn) {$\cdots$};
\node[smalllab] at (xn) {$x_n$};

\foreach \pt/\lbl in {p1/{\phi_{q,1}},p2/{\phi_{q,2}},p3/{\phi_{q,3}},pn/{\phi_{q,n}}}{
    \draw[accent, line width=0.75pt]
        ($(\pt)+(-0.24,-0.05)$)
        .. controls +(0.10,0.16) and +(-0.14,-0.10) ..
        ($(\pt)+(0.03,0.18)$)
        .. controls +(0.10,0.12) and +(-0.12,-0.08) ..
        ($(\pt)+(0.24,0.00)$);
    \node[smalllab, accent] at ($(\pt)+(0,0.36)$) {$\lbl$};
}

\draw[plainarrow] (x1) -- (p1);
\draw[plainarrow] (x2) -- (p2);
\draw[plainarrow] (x3) -- (p3);
\draw[plainarrow] (xn) -- (pn);
\draw[plainarrow] (p1) -- ($(LAgg.south)+(-1.02,0.02)$);
\draw[plainarrow] (p2) -- ($(LAgg.south)+(-0.33,0.02)$);
\draw[plainarrow] (p3) -- ($(LAgg.south)+(0.33,0.02)$);
\draw[plainarrow] (pn) -- ($(LAgg.south)+(1.02,0.02)$);

% Center column: HMF
\node[box, text width=4.80cm, font=\scriptsize] (outC) at (\xC,\yOut)
    {$\begin{aligned}
        H_{\mathrm{MF}}(\beta) &= H_Q-\frac{1}{\beta}
        {\color{accent}\Phi}\!\big(-\beta\,\mathrm{ad}_{H_Q}\big)\Delta \\
        &\quad + O(\Delta^2)
    \end{aligned}$};

\coordinate (CG) at (7.70,\yOuter);
\draw[muted, thin] ($(CG)+(-0.70,-0.56)$) -- ($(CG)+(0.86,-0.56)$);
\draw[muted, thin] ($(CG)+(0,-0.76)$) -- ($(CG)+(0,0.96)$);
\draw[accent, line width=1.0pt]
    plot[smooth] coordinates {
        ($(CG)+(-0.56,-0.44)$)
        ($(CG)+(-0.34,-0.27)$)
        ($(CG)+(-0.12,-0.05)$)
        ($(CG)+(0.08,0.18)$)
        ($(CG)+(0.30,0.47)$)
        ($(CG)+(0.52,0.78)$)
        ($(CG)+(0.74,1.10)$)
    };
\node[phimark] at (8.28,\yOuter+0.76) {$\Phi$};
\node[smalllab, text=muted] at (8.92,\yOuter+0.92) {fixed};
\node[align=left, anchor=west, font=\scriptsize] (PhiText) at (8.55,\yOuter+0.05)
    {$\textcolor{accent}{\Phi(x)}=\dfrac{x}{1-e^{-x}}$\\[2pt]
     $\textcolor{accent}{\Phi(x)}=B_0 + B_1(-x) + \dfrac{B_2}{2!}x^2$};
\draw[plainarrow] ($(CG)+(0.10,0.78)$) -- (outC.south);

\draw[thin, rounded corners=2pt] (6.35,4.28) rectangle (10.85,5.86);
\node[align=left, text width=2.95cm] at (8.08,\yAgg)
    {$\displaystyle \Delta(\beta)=\sum_{n,m} C^>_{nm}(\beta)\,f_n f_m$};
\foreach \i/\x in {0/9.90,1/10.12,2/10.34}{
    \foreach \j/\y/\shade in {
        0/5.42/20,
        1/5.20/35,
        2/4.98/55
    }{
        \pgfmathtruncatemacro{\mix}{\shade + 10*\i}
    \fill[accent!\mix, draw=white] (\x,\y) rectangle ++(0.18,0.18);
    }
}
\draw[thin] (9.90,4.98) rectangle (10.48,5.56);
\node[smalllab] at (10.19,4.72) {$C^>_{nm}$};
\draw[plainarrow] (8.90,5.86) -- ($(CG)+(0,-0.66)$);

\node[orbitnode] (f0) at (\xC,0.88) {$f_0=f$};
\node[orbitnode] (f1) at (\xC,1.86) {$f_1$};
\node[orbitnode] (f2) at (\xC,2.84) {$f_2$};
\node[smalllab, text=muted] at (\xC,3.64) {$\vdots$};
\draw[plainarrow] (f0) -- (f1);
\draw[plainarrow] (f1) -- (f2);
\draw[plainarrow] (f2) -- (\xC,3.35);
\node[smalllab, rotate=90, text=muted] at (9.64,2.36) {$\mathrm{ad}_{H_Q}$};
\draw[plainarrow] (\xC,3.88) -- (8.90,4.28);

% Right column: exact qubit closure
\node[box, text width=3.05cm] (outR) at (\xR,\yOut)
    {$H_{\mathrm{MF}}^{\mathrm{qb}}(\beta)$\\[-1pt]
    \scriptsize closed Bloch-plane generator};

\node[orbitnode, minimum width=0.75cm, minimum height=0.75cm] (Mnode) at (\xR,\yOuter) {$M$};
\draw[plainarrow] ($(Mnode)+(0.82,0)$) arc[start angle=0,end angle=330,radius=0.82];
\node[smalllab, anchor=west] at (\xR+0.72,\yOuter+0.18) {$M^2=\chi^2\mathbb I$};
\draw[plainarrow] (Mnode.north) -- (outR.south);

\coordinate (qsq) at (\xR,\yAgg);
\node[orbitnode, minimum width=0.78cm] (Id) at (\xR-0.56,\yAgg+0.50) {$\mathbb I$};
\node[orbitnode, minimum width=0.78cm] (Sz) at (\xR+0.56,\yAgg+0.50) {$\sigma_z$};
\node[orbitnode, minimum width=0.78cm] (Sx) at (\xR-0.56,\yAgg-0.50) {$\sigma_x$};
\node[orbitnode, minimum width=0.78cm] (Sy) at (\xR+0.56,\yAgg-0.50) {$\sigma_y$};
\draw[thin] (Id) -- (Sz) -- (Sy) -- (Sx) -- cycle;
\node[smalllab, align=center, text width=3.00cm] at (\xR,\yAgg-1.08)
    {$\Delta \in \mathrm{span}\{\mathbb I,\sigma_x,\sigma_y,\sigma_z\}$};
\draw[plainarrow] (\xR,\yAgg+0.94) -- (Mnode.south);

\node[orbitnode] (q0) at (\xR,0.82) {$f_0$};
\node[orbitnode, text width=1.48cm, minimum width=1.55cm] (qe) at (\xR,1.72)
    {$f_{2k}\!\sim\!\sigma_x$};
\node[orbitnode, text width=1.48cm, minimum width=1.55cm] (qo) at (\xR,2.92)
    {$f_{2k+1}\!\sim\!\sigma_y$};
\draw[plainarrow] (q0) -- (qe);
\draw[plainarrow] (qe) to[bend left=40] node[right=2pt, smalllab] {$\mathrm{ad}_{H_Q}$} (qo);
\draw[plainarrow] (qo) to[bend left=40] (qe);
\draw[plainarrow] (qo.north) -- (\xR,4.28);

\end{tikzpicture}
}
    \caption{Three-layer compositional correspondence between
    Kolmogorov--Arnold representations, the exact Gaussian Hamiltonian of mean
    force, and the spin-boson qubit. Left: a KAN builds a target from
    univariate edge functions $\phi_{q,p}(x_p)$, additive aggregation, and
    outer univariate maps $\Phi_q$. Center: the HMF construction has the same
    layer structure, with adjoint orbit
    $f_n=\mathrm{ad}_{H_Q}^n(f)$, bilinear aggregation
    $\Delta(\beta)=\sum_{n,m} C^>_{nm}(\beta)f_n f_m$, and outer composition by
    the Todd function $\Phi(x)=x/(1-e^{-x})$ fixed analytically by the BCH
    expansion. Right: for the qubit, Conditions (C1)--(C3) close the adjoint
    orbit, bilinear sector, and BCH recombination exactly. The shared
    accent-coloured $\Phi$ highlights the visual rhyme between learned and
    thermodynamically fixed outer composition.}
    \label{fig:kan_hmf_closure}
\end{figure*}

This is more than a family resemblance. The Gaussian HMF construction already
provides the three ingredients usually separated in a compact compositional
model: basis atoms, an aggregation tensor, and a nonlinear composition rule.
The adjoint orbit is a discrete univariate basis in operator space: one
generator, one seed, one orbit index $n$. The matrix $C^>_{nm}(\beta)$ is the
aggregation tensor, but here it is fixed analytically by the bath kernel rather
than learned from data. The Todd/Bernoulli superoperator
\(\Phi(-\beta\,\mathrm{ad}_{H_Q})\), together with the higher BCH nests,
supplies the nonlinear composition rule. Thermodynamics therefore provides both
the architecture and the weights.

Once written in this form, representability failure becomes a precise statement
about model-class insufficiency. If the adjoint orbit or its quadratic and BCH
closures escape $\mathcal A$, the exact reduced generator still exists but no
finite ansatz in that class can carry it without projection. For non-quadratic
continuous-variable potentials the orbit typically expands into the universal
enveloping algebra, so the thermodynamic target outgrows any fixed finite
operator family. That is the same structural obstruction that motivates compact
decompositions in KANs: the universal object exists, but the chosen finite
representation need not be rich enough to realize it.

In that more modest sense the discussion also touches mechanistic
interpretability. Once the layers are explicit, one can inspect which channels
carry the reduced generator and which ones force basis growth. The object being
interpreted here, however, is not a trained network but an exact reduced
equilibrium operator. We now exhibit a case where all three compositional
layers close exactly.
